% ============================================================================
% ARBEITSBLATT: Reibungskräfte (komplett - Doppelstunde)
% Physik Klasse 7 - Kräfte
% Stand: 04.02.2026
% ============================================================================

\documentclass[11pt,a4paper]{article}

\usepackage[utf8]{inputenc}
\usepackage[T1]{fontenc}
\usepackage[ngerman]{babel}
\usepackage{geometry}
\usepackage{helvet}
\usepackage{enumitem}
\usepackage{fancyhdr}
\usepackage{xcolor}
\usepackage{tcolorbox}
\usepackage{amssymb}
\usepackage{siunitx}
\usepackage{qrcode}
\usepackage{array}
\usepackage{tabularx}
\usepackage{textcomp}
\usepackage[hidelinks]{hyperref}

\sisetup{locale=DE, per-mode=power, inter-unit-product=\ensuremath{\cdot}}
\renewcommand{\familydefault}{\sfdefault}

\geometry{left=1.4cm, right=1.4cm, top=1.3cm, bottom=1.0cm}

% Fachfarbe Physik
\definecolor{physik}{HTML}{2c5aa0}
\colorlet{fachfarbe}{physik}
\definecolor{protokollrand}{HTML}{2a7a4b}
\definecolor{merkerand}{HTML}{b35c00}

% Kopfzeile
\pagestyle{fancy}
\fancyhf{}
\fancyhead[L]{\small\textbf{Physik Klasse 7} | Reibungskräfte}
\fancyhead[R]{\small Arbeitsblatt}
\renewcommand{\headrulewidth}{0.4pt}
\setlength{\headheight}{14pt}

% Boxen
\tcbuselibrary{skins}

\newtcolorbox{merkkasten}[1][]{
    colback=white, colframe=fachfarbe,
    coltitle=black, colbacktitle=white,
    fonttitle=\bfseries\small,
    boxrule=1pt, arc=0pt,
    left=6pt, right=6pt, top=4pt, bottom=4pt,
    #1
}

\newtcolorbox{protokollbox}[1][]{
    colback=white, colframe=protokollrand,
    coltitle=black, colbacktitle=white,
    fonttitle=\bfseries\small,
    boxrule=1pt, arc=0pt,
    left=6pt, right=6pt, top=4pt, bottom=4pt,
    #1
}

\newtcolorbox{merkebox}[1][]{
    colback=white, colframe=merkerand,
    coltitle=black, colbacktitle=white,
    fonttitle=\bfseries\small,
    boxrule=1pt, arc=0pt,
    left=6pt, right=6pt, top=4pt, bottom=4pt,
    #1
}

\newcommand{\aufgabe}[1]{\noindent\textbf{\textcolor{fachfarbe}{#1}}}
\newcommand{\luecke}[1]{\underline{\hspace{#1}}}
\newcommand{\stern}[1]{\textcolor{merkerand}{\textbf{#1}}}

\setlength{\parindent}{0pt}
\setlength{\parskip}{0pt}

\begin{document}

% ============================================================================
% TITEL
% ============================================================================
\begin{minipage}[t]{0.80\textwidth}
    \vspace{0pt}
    {\Large\bfseries Reibungskräfte}\\[0.2em]
    {\small Name: \luecke{4.5cm} \hfill Datum: \luecke{2.5cm}}
\end{minipage}
\hfill
\begin{minipage}[t]{0.15\textwidth}
    \vspace{0pt}
    \raggedleft
    \qrcode[height=1.4cm]{https://devbox42.github.io/sciencesim/physik/kl07/02-kraefte/05-reibungskraefte/LP-05-komplett.html}
\end{minipage}
\vspace{0.1em}
\hrule
\vspace{0.4em}

% ============================================================================
% K1: INFOKASTEN
% ============================================================================
\begin{merkkasten}[title=K1: Was ist Reibung?]
\small
Reibung ist eine \textbf{Kraft}, die einer Bewegung entgegenwirkt. Sie entsteht, wenn zwei Oberflächen aneinander vorbeigleiten oder gleiten wollen.

\vspace{0.3em}
\textbf{Beispiele im Alltag:} Bremsen, Schuhe auf dem Boden, Streichholz anzünden, ...
\end{merkkasten}

\vspace{0.3em}

% ============================================================================
% B1: BEOBACHTUNG HAFT/GLEIT
% ============================================================================
\aufgabe{B1: Beobachtung -- Klotz anschieben}

\vspace{0.2em}
Beobachte die Animation und ergänze:

\vspace{0.2em}
Um den Klotz in Bewegung zu setzen, muss die Kraft erst \luecke{4cm}.

Sobald der Klotz gleitet, \luecke{3cm} die benötigte Kraft.

\vspace{0.4em}

% ============================================================================
% T1: MESSWERTE
% ============================================================================
\begin{protokollbox}[title=T1: Messwerte -- Drei Reibungsarten]
\small
Miss in der Simulation und trage die Werte ein:

\vspace{0.3em}
\renewcommand{\arraystretch}{1.4}
\begin{tabular}{|p{4cm}|p{3cm}|p{5cm}|}
\hline
\textbf{Reibungsart} & \textbf{Messwert} & \textbf{Beschreibung} \\
\hline
Haftreibung (Losbrechkraft) & \luecke{1.5cm} N & Kraft zum Lösen \\
\hline
Gleitreibung & \luecke{1.5cm} N & Kraft während Bewegung \\
\hline
Rollreibung & \luecke{1.5cm} N & Kraft beim Rollen \\
\hline
\end{tabular}
\end{protokollbox}

\vspace{0.3em}

% ============================================================================
% B2/K2: ERKENNTNIS
% ============================================================================
\aufgabe{B2: Erkenntnis -- Die Reihenfolge}

\vspace{0.2em}
Setze die richtigen Zeichen ein ($>$ oder $<$):

\vspace{0.2em}
\begin{center}
{\large $F_{\text{Haft}}$ \luecke{1cm} $F_{\text{Gleit}}$ \luecke{1cm} $F_{\text{Roll}}$}
\end{center}

\vspace{0.3em}
\begin{merkebox}[title=K2: Merksatz -- Die drei Reibungsarten]
\small
Die \luecke{3cm} ist am größten, die \luecke{3cm} am kleinsten.

\vspace{0.2em}
$\rightarrow$ Deshalb ist das Erfinden des \luecke{2cm} so wichtig: Rollreibung statt Gleitreibung!
\end{merkebox}

\vspace{0.3em}

% ============================================================================
% B3/B4: EINFLUSSFAKTOREN
% ============================================================================
\aufgabe{B3/B4: Einflussfaktoren der Reibung}

\vspace{0.2em}
Experimentiere mit der Simulation und ergänze:

\vspace{0.2em}
\begin{tabular}{ll}
B3: & Mehr Gewichtskraft = \luecke{2cm} Reibung \\[0.3em]
B4: & Rauere Oberfläche = \luecke{2cm} Reibung \\
\end{tabular}

\vspace{0.4em}

% ============================================================================
% T2: ERWUENSCHT/UNERWUENSCHT
% ============================================================================
\begin{protokollbox}[title=T2: Erwünschte und unerwünschte Reibung]
\small
\renewcommand{\arraystretch}{1.3}
\begin{tabular}{|p{5.8cm}|p{5.8cm}|}
\hline
\textbf{Erwünscht (nützlich)} & \textbf{Unerwünscht (störend)} \\
\hline
1. \luecke{4cm} & 1. \luecke{4cm} \\[0.5em]
2. \luecke{4cm} & 2. \luecke{4cm} \\[0.5em]
3. \luecke{4cm} & 3. \luecke{4cm} \\
\hline
\end{tabular}
\end{protokollbox}

\vspace{0.2em}

\aufgabe{A1: Fahrrad} Wo ist beim Fahrrad Reibung erwünscht, wo unerwünscht? Nenne je ein Beispiel.

\vspace{0.2em}
Erwünscht: \luecke{8cm}

Unerwünscht: \luecke{8cm}

% ============================================================================
% SEITE 2
% ============================================================================
\newpage

% ============================================================================
% K3: FORMEL + DREIECK
% ============================================================================
\begin{merkebox}[title=K3: Die Reibungsformel \stern{$\star\star$}]
\small
\textbf{Die Formel:} \quad {\large $F_\text{R} = \mu \cdot F_\text{N}$}

\vspace{0.3em}
\renewcommand{\arraystretch}{1.2}
\begin{tabular}{ll}
$F_\text{R}$ & = Reibungskraft in \luecke{1cm} \\
$\mu$ & = Reibungszahl (``mü''), \luecke{2cm} Einheit \\
$F_\text{N}$ & = Normalkraft in \luecke{1cm} (auf ebener Fläche $\approx$ Gewichtskraft) \\
\end{tabular}

\vspace{0.4em}
\textbf{Das Formeldreieck:} Zeichne es ab und notiere alle drei Formeln!

\vspace{0.3em}
\begin{minipage}{0.35\textwidth}
\centering
{\scriptsize (Zeichnung hier)}
\vspace{2.5cm}
\end{minipage}
\hfill
\begin{minipage}{0.58\textwidth}
\renewcommand{\arraystretch}{1.5}
\begin{tabular}{|p{2cm}|p{5cm}|}
\hline
\textbf{Gesucht} & \textbf{Formel} \\
\hline
$F_\text{R}$ & $F_\text{R} =$ \luecke{3cm} \\
\hline
$\mu$ & $\mu =$ \luecke{3cm} \\
\hline
$F_\text{N}$ & $F_N =$ \luecke{3cm} \\
\hline
\end{tabular}
\end{minipage}

\vspace{0.3em}
\textbf{Reibungszahl-Tabelle:}
\begin{tabular}{lll}
Gummi -- Asphalt: $\mu = 0{,}8$ & Holz -- Holz: $\mu = 0{,}5$ & Stahl -- Eis: $\mu = 0{,}03$
\end{tabular}
\end{merkebox}

\vspace{0.3em}

% ============================================================================
% RECHENAUFGABEN
% ============================================================================
\aufgabe{Rechenaufgaben} {\small Zeige den vollständigen Rechenweg: Gegeben, Gesucht, Formel, Einsetzen, Ergebnis!}

\vspace{0.3em}
\textbf{\stern{$\star$} R1:} Ein Klotz liegt auf dem Tisch. $\mu = 0{,}4$, $F_N = \SI{20}{N}$. Berechne $F_\text{R}$.

\vspace{1.5cm}

\textbf{\stern{$\star\star$} R2:} Autoreifen auf trockener Straße: $\mu = 0{,}8$, $F_N = \SI{5000}{N}$. Berechne $F_\text{R}$.

\vspace{1.5cm}

\textbf{\stern{$\star\star\star$} R3:} Die Reibungskraft beträgt $F_\text{R} = \SI{12}{N}$, die Normalkraft $F_N = \SI{40}{N}$. Berechne $\mu$.

\vspace{1.5cm}

% ============================================================================
% A2: BREMSWEG
% ============================================================================
\begin{protokollbox}[title=A2: Anwendung -- Bremsweg]
\small
Beobachte die Bremsweg-Animation und beantworte:

\vspace{0.3em}
a) Welche Reibungsart ist beim Bremsen wichtig? \luecke{3cm}

\vspace{0.3em}
b) Auf welcher Straße ist der Bremsweg am längsten? \luecke{3cm}

\vspace{0.3em}
c) Erkläre in einem Satz, warum: \luecke{7cm}

\luecke{12cm}
\end{protokollbox}

\vspace{0.4em}

% ============================================================================
% SELBSTCHECK
% ============================================================================
\begin{merkkasten}
\textbf{Selbstcheck: Das kann ich jetzt!} \hfill {\scriptsize \textcircled{\scriptsize 1} = unsicher \quad \textcircled{\scriptsize 2} = geht so \quad \textcircled{\scriptsize 3} = sicher}

\vspace{0.2em}
\small
\renewcommand{\arraystretch}{1.2}
\begin{tabular}{p{10cm}ccc}
Ich kann die drei Reibungsarten nennen und nach Größe ordnen. & $\square$ & $\square$ & $\square$ \\
Ich kann erklären, wovon die Reibungskraft abhängt. & $\square$ & $\square$ & $\square$ \\
Ich kann erwünschte und unerwünschte Reibung unterscheiden. & $\square$ & $\square$ & $\square$ \\
\stern{$\star\star$} Ich kann die Formel $F_\text{R} = \mu \cdot F_\text{N}$ anwenden. & $\square$ & $\square$ & $\square$ \\
\stern{$\star\star\star$} Ich kann die Formel umstellen und $\mu$ berechnen. & $\square$ & $\square$ & $\square$ \\
\end{tabular}
\end{merkkasten}

\end{document}
